%=====================================================================
\chapter{Heaps and Search Trees, Revisited for Analysis}\label{ch:structures}
%=====================================================================
\locator{3}{Part III --- Sorting, Selection and the Structures Beneath}

\begin{outcomes}
By the end of this chapter you will be able to:
\begin{tight}
  \item \textbf{Explain} why a binary heap's height is exactly
        $\lfloor \lg n \rfloor$ and why that cannot fail.
  \item \textbf{Prove} that \algname{Build-Heap} runs in $\bigTh{n}$, not
        $\bigTh{n \lg n}$, and \textbf{identify} the step where the naive
        analysis is loose.
  \item \textbf{Contrast} a heap's structural guarantee with a binary search
        tree's dependence on insertion order \emph{(CLO-2)}.
  \item \textbf{State} the expected height of a randomly built BST and
        \textbf{assess} measured data against it.
  \item \textbf{Decide} between a heap, a balanced tree and a sorted array for
        a stated workload.
\end{tight}
\end{outcomes}

\begin{prereq}
\emph{Data Structures} built all of these; this chapter assumes you can recall
what a heap and a BST \emph{are}. \chref{recurrences} for the summation in
Section~12.2, and \chref{cases} for the best/worst-case vocabulary.
\end{prereq}

\begin{keyterms}
binary heap $\bullet$ complete binary tree $\bullet$ heap property $\bullet$
sift down $\bullet$ build-heap $\bullet$ height $\bullet$ binary search tree
$\bullet$ degenerate tree $\bullet$ internal path length $\bullet$ randomly
built BST $\bullet$ balanced tree $\bullet$ rotation
\end{keyterms}

\begin{alertbox}[What this chapter is, and is not]
It is \textbf{not} a reimplementation of heaps and search trees. Appendix~F
records the boundary: \emph{Data Structures} owns the structures and their
operations; this course owns their analysis.

\smallskip
So the question here is never ``how do you insert into a heap''. It is
\emph{why is the height what it is, what guarantees it, and what does a broken
guarantee cost?} --- and the answers divide these two structures sharply.
\end{alertbox}

\section{A Heap's Height Is a Guarantee}

\begin{definitionbox}[Binary Heap]
A \term{binary heap} is a \term{complete binary tree} --- every level full
except possibly the last, which is filled from the left --- satisfying the
\term{heap property}: every node's key is $\ge$ (max-heap) or $\le$ (min-heap)
its children's.

\smallskip
It is stored in an array, with the children of index $i$ at $2i+1$ and $2i+2$.
\end{definitionbox}

\begin{theorem}[Heap height]\label{thm:heapheight}
A binary heap on $n$ elements has height exactly $\lfloor \lg n \rfloor$.
\end{theorem}

\begin{proof}
A complete binary tree of height $h$ has at least $2^{h}$ nodes --- levels $0$
to $h-1$ are full, contributing $2^{h}-1$, plus at least one node on level $h$
--- and at most $2^{h+1}-1$, when every level is full. So
\[ 2^{h} \le n \le 2^{h+1}-1 \quad\Longrightarrow\quad
   h \le \lg n < h+1 \quad\Longrightarrow\quad h = \lfloor \lg n \rfloor. \qedhere \]
\end{proof}

\begin{conceptbox}[Why this cannot go wrong]
Theorem~\ref{thm:heapheight} has no hypothesis about the data. Compare that
with everything else in this book.

\smallskip
The height is $\lfloor \lg n \rfloor$ because the tree is \emph{complete}, and
it is complete because of the \emph{array layout}: node $i$'s children are at
$2i+1$ and $2i+2$, and a new element goes at index $n$. There is no insertion
order that produces an unbalanced heap, because there is no representation in
which an unbalanced heap could be stored.

\smallskip
\textbf{The structure is enforced by the addressing scheme, not maintained by
an algorithm.} That is why heaps need no rebalancing, no rotations, no
colouring and no amortised argument --- and it is the sharpest contrast in this
chapter with Section~12.3.
\end{conceptbox}

\begin{worked}[Measured, Because It Costs Nothing To Check]
\texttt{code/ch12\_structures.cpp} walks from the last node to the root:

\rowcolors{2}{white}{lightbg}
\begin{center}
\begin{tabular}{@{}rrrl@{}}
\toprule
\rowcolor{primary!15}
$n$ & height & $\lfloor \lg n \rfloor$ & \\
\midrule
      1{,}000 &  9 &  9 & equal \\
    100{,}000 & 16 & 16 & equal \\
 10{,}000{,}000 & 23 & 23 & equal \\
100{,}000{,}000 & 26 & 26 & equal \\
\bottomrule
\end{tabular}
\end{center}

\smallskip
Exact at every size, as a theorem with no hypotheses should be. A hundred
million elements are 26 levels deep.
\end{worked}

\section{Build-Heap Is Linear, and the Obvious Analysis Says Otherwise}

\dsfig{\aoaalg{\algname{Build-Heap}$(A, n)$}{
\For{$i \gets \lfloor n/2 \rfloor - 1$ \textbf{downto} $0$}
  \State \algname{Sift-Down}$(A, i, n)$
    \Comment{nodes below $\lfloor n/2 \rfloor$ are leaves}
\EndFor
}}{\algname{Build-Heap}. The loop runs $n/2$ times and \algname{Sift-Down}
costs $\bigO{\lg n}$, so the obvious bound is $\bigO{n \lg n}$. The obvious
bound is loose, and Theorem~\ref{thm:buildheap} says by how much.}{buildheap}

\begin{theorem}[Build-Heap]\label{thm:buildheap}
\algname{Build-Heap} runs in $\bigTh{n}$ time.
\end{theorem}

\begin{proof}
\algname{Sift-Down} on a node at height $j$ --- counting height upward from the
leaves --- costs $\bigO{j}$, since it travels at most $j$ levels down.

\smallskip
The key observation is that \emph{most nodes are near the bottom}. In a heap of
$n$ elements there are at most $\lceil n/2^{j+1} \rceil$ nodes of height $j$.
So the total cost is
\[ \sum_{j=0}^{\lfloor \lg n \rfloor} \Bigl\lceil \frac{n}{2^{j+1}} \Bigr\rceil
   \bigO{j}
   \;=\; \bigO{\,n \sum_{j=0}^{\lfloor \lg n \rfloor} \frac{j}{2^{j}}\,}. \]
The series $\sum_{j\ge0} j/2^{j}$ converges --- its value is exactly 2, from
$\sum_{j\ge0} jx^{j} = x/(1-x)^{2}$ at $x = 1/2$. Hence the total is
$\bigO{2n} = \bigO{n}$, and it is $\bigOm{n}$ because every element is
examined, giving $\bigTh{n}$.
\end{proof}

\begin{conceptbox}[Where the naive analysis was loose]
``$n/2$ calls, each $\bigO{\lg n}$, therefore $\bigO{n \lg n}$'' is valid ---
it is a correct upper bound. It is simply not tight, and the reason is worth
internalising.

\smallskip
It charges \emph{every} node the cost of the \emph{worst} node. But half the
nodes are leaves and cost nothing, a quarter are one level up and cost 1, an
eighth cost 2. The expensive nodes are the rare ones, and the sum
$\sum j/2^{j}$ converges precisely because the cost grows linearly while the
population falls geometrically.

\smallskip
\textbf{This pattern --- a bound that is loose because it multiplies a worst
case by a count --- recurs throughout the book}, and \chref{amortised} is
entirely devoted to it. Heapsort, by contrast, really is $\bigTh{n\lg n}$: it
calls \algname{Sift-Down} $n$ times from the \emph{root}, where the height
really is $\lg n$ every time.
\end{conceptbox}

\section{A Search Tree's Height Is a Hope}

\begin{alertbox}[The same keys, the same code, two different structures]
A BST has no array layout to enforce its shape. Its height is decided entirely
by the \emph{order the keys arrive in} --- which is \chref{cases}'s lesson,
applied to a data structure rather than to a sort.

\smallskip
Insert $1, 2, 3, \ldots, n$ in that order and every key is larger than
everything before it, so every key goes down the right spine. The result has
height $n$: it is a linked list that costs more to traverse than a linked list,
because the pointers are scattered.
\end{alertbox}

\begin{worked}[Measured: Random Keys Against Sorted Keys]
\rowcolors{2}{white}{lightbg}
\begin{center}
\begin{tabular}{@{}rrrrrr@{}}
\toprule
\rowcolor{primary!15}
$n$ & $\lg n$ & random $h$ & $h/\lg n$ & sorted $h$ & sorted $h/n$ \\
\midrule
  1{,}000 & 10.0 & 21 & 2.11 &   1{,}000 & 1.000 \\
  8{,}000 & 13.0 & 27 & 2.08 &   8{,}000 & 1.000 \\
 64{,}000 & 16.0 & 37 & 2.32 &  64{,}000 & 1.000 \\
256{,}000 & 18.0 & 44 & 2.45 & 256{,}000 & 1.000 \\
\bottomrule
\end{tabular}
\end{center}

\smallskip
\textbf{The sorted column is $n$ exactly} --- $h/n = 1.000$ at every size. Not
``about $n$''; exactly $n$.

\smallskip
\textbf{The random column is about $2.3 \lg n$}, which is a tree. At
$n = 256{,}000$: height 44 against height 256{,}000. A factor of 5{,}800, from
nothing but the order the same keys were handed over in.
\end{worked}

\begin{conceptbox}[The constants for a randomly built BST]
Two classical results, both asymptotic:

\begin{tight}
  \item \textbf{Expected height} $\sim 4.311 \ln n - 1.953 \ln\ln n + \bigO{1}$
        (Reed, 2003). The leading term is $2.988 \lg n$.
  \item \textbf{Expected node depth} $\sim 2\ln n - 3 + \bigO{1}$. The leading
        term is $1.386 \lg n$, and this is the average cost of a successful
        search.
\end{tight}

\smallskip
Note that ``randomly built'' means the keys were \emph{inserted} in random
order. It does not mean the keys are random --- sorted keys inserted in sorted
order is the disaster above, and the keys there were as random as you like.
\end{conceptbox}

\begin{worked}[The Constants, Measured Honestly]
Mean of 25 independently shuffled trees. Both measured columns are divided by
$\lg n$.

\rowcolors{2}{white}{lightbg}
\begin{center}
\begin{tabular}{@{}rrrrrrr@{}}
\toprule
\rowcolor{primary!15}
$n$ & height & corrected & asympt. & mean depth & corrected & asympt. \\
\midrule
  2{,}000 & 2.324 & 2.627 & 2.988 & 1.231 & 1.113 & 1.386 \\
  8{,}000 & 2.311 & 2.657 & 2.988 & 1.242 & 1.155 & 1.386 \\
 32{,}000 & 2.459 & 2.683 & 2.988 & 1.276 & 1.186 & 1.386 \\
128{,}000 & 2.461 & 2.704 & 2.988 & 1.278 & 1.209 & 1.386 \\
\bottomrule
\end{tabular}
\end{center}

\smallskip
\textbf{Quoting only the asymptotic constants would have made the measurement
look wrong.} Measured height$/\lg n$ is $2.46$ against an asymptote of $2.99$
--- an 18\% shortfall, which looks like a failed prediction until you include
the $-1.953\ln\ln n$ term. At $n = 128{,}000$, $\ln\ln n = 2.46$, and that term
alone subtracts $0.28$ from the ratio.

\smallskip
\textbf{The mean-depth column is the cleaner story.} Measured $1.231 \to 1.278$,
rising steadily towards the asymptote $1.386$, and bracketed above the
corrected form $(2\ln n - 3)/\lg n$. It is converging, slowly, from below ---
and ``slowly'' here means like $1/\lg n$, so reaching within 1\% of the
asymptote would need $n$ around $10^{30}$.

\smallskip
\textbf{The general lesson is about reading asymptotic constants.} An
asymptotic constant is a limit, and the lower-order terms it discards can
dominate at every size you will ever run. When a measurement misses a
theoretical constant by 20\%, the first question is not ``is the theory wrong''
but ``what did the $\bigO{1}$ swallow''.
\end{worked}

\begin{worked}[What the Height Costs, in Seconds]
100{,}000 successful searches against each tree:

\rowcolors{2}{white}{lightbg}
\begin{center}
\begin{tabular}{@{}rrrr@{}}
\toprule
\rowcolor{primary!15}
$n$ & random BST (ms) & sorted BST (ms) & ratio \\
\midrule
 2{,}000 &  6.65 &   532.05 &  80.0 \\
 8{,}000 & 10.02 & 1{,}886.29 & 188.3 \\
16{,}000 &  9.40 & 3{,}790.47 & 403.3 \\
32{,}000 & 11.72 & 7{,}618.81 & 650.1 \\
\bottomrule
\end{tabular}
\end{center}

\smallskip
\textbf{The random tree's time barely moves} --- 6.6 ms to 11.7 ms while $n$
grows 16-fold --- which is what $\lg n$ looks like on a clock.

\smallskip
\textbf{The degenerate tree's time doubles with $n$}, and the ratio reaches
650$\times$ at only 32{,}000 elements.

\smallskip
\textbf{And it is worse than the $n/\lg n$ ratio predicts.} At $n = 32{,}000$,
$n/(2.4\lg n) \approx 890$ steps against 44, a ratio of about 20 in
\emph{steps} --- yet the measured time ratio is 650. The extra factor is memory:
each step in the degenerate tree is a dependent pointer dereference to an
unpredictable address, so the processor can neither prefetch it nor overlap it
with the next. \chref{role}'s second lie --- that all memory costs the same ---
charging interest again.
\end{worked}

\section{Choosing}

\begin{conceptbox}[Three structures, one table]
\rowcolors{2}{white}{lightbg}
\begin{longtable}{@{}L{2.8cm}L{2.4cm}L{2.4cm}L{2.4cm}L{3.3cm}@{}}
\toprule
\rowcolor{primary!15}
\textbf{Operation} & \textbf{Sorted array} & \textbf{Binary heap} &
\textbf{Unbalanced BST} & \textbf{Balanced BST} \\
\midrule
\endhead
find min/max & $\bigTh{1}$ & $\bigTh{1}$ & $\bigTh{h}$ & $\bigTh{\lg n}$ \\
search & $\bigTh{\lg n}$ & $\bigTh{n}$ & $\bigTh{h}$ & $\bigTh{\lg n}$ \\
insert & $\bigTh{n}$ & $\bigTh{\lg n}$ & $\bigTh{h}$ & $\bigTh{\lg n}$ \\
delete min & $\bigTh{n}$ & $\bigTh{\lg n}$ & $\bigTh{h}$ & $\bigTh{\lg n}$ \\
in-order walk & $\bigTh{n}$ & $\bigTh{n\lg n}$ & $\bigTh{n}$ & $\bigTh{n}$ \\
\textbf{worst-case $h$} & --- & $\lfloor\lg n\rfloor$ \emph{guaranteed} &
  $n$ & $\bigO{\lg n}$ \emph{guaranteed} \\
\bottomrule
\end{longtable}

\smallskip
\textbf{A heap is not a search structure.} It gives $\bigTh{1}$ access to the
extremum and $\bigTh{n}$ search for anything else --- the heap property orders
each node against its children only, never its siblings. Use it for priority
queues (\chref{mst}, \chref{shortestpaths}) and nothing else.

\smallskip
\textbf{An unbalanced BST is a gamble on insertion order}, and the measured
penalty for losing is 650$\times$. Balanced variants --- AVL, red--black,
treaps --- restore the guarantee by rotating on insertion, at a constant-factor
cost. \emph{Data Structures} covers the rotations; what this course adds is
that the guarantee is worth paying for, and by how much.
\end{conceptbox}

\begin{pitfall}
\begin{tight}
  \item \textbf{Assuming \algname{Build-Heap} is $\bigO{n\lg n}$.} It is
        $\bigTh{n}$; the loose analysis charges every node the root's cost.
  \item \textbf{Searching a heap.} It is $\bigTh{n}$. Heaps order parents
        against children, not siblings.
  \item \textbf{Expecting a BST to be balanced.} Its height is decided by the
        insertion order, and sorted input gives height exactly $n$.
  \item \textbf{Testing a BST only with shuffled keys.} The disaster case is
        \emph{ordered} input, which is what real data often is.
  \item \textbf{Comparing a measured constant with a bare asymptote.} The
        height ratio was 2.46 against an asymptote of 2.99 and the theory was
        right; the $\ln\ln n$ term accounted for the difference.
  \item \textbf{Predicting slowdown from step counts alone.} The measured 650×
        came from about 20× in steps and the rest from unpredictable memory
        access.
  \item \textbf{Reaching for a balanced tree when a sorted array would do.} If
        the data is static, binary search on an array is $\bigTh{\lg n}$ with a
        far better constant and no pointers at all.
\end{tight}
\end{pitfall}

\begin{lab}[Height, Measured]
Reproduces all four tables. Expect twenty-five minutes; part D is slow by
design --- it is timing a deliberately degenerate structure.

\begin{lstlisting}[language=C++]
// ch12_structures.cpp -- heaps and search trees, analysed not rebuilt.
// Build: cl /O2 /EHsc /std:c++17 ch12_structures.cpp (see Appendix A)

// --- 1. A BST as three parallel arrays -----------------------------
// No pointers and no allocation: left[i], right[i] and key[i]. The
// shape is identical to a pointer tree, and the measurement is about
// shape, not about allocators.
    void insert(int k) {
        if (root < 0) { root = make(k); return; }
        int cur = root;
        for (;;) {
            if (k < key[cur]) {
                if (left[cur] < 0) { left[cur] = make(k); return; }
                cur = left[cur];
            } else {
                if (right[cur] < 0) { right[cur] = make(k); return; }
                cur = right[cur];
            }
        }
    }
// Iterative height, so a degenerate tree does not overflow the stack
// the way Chapter 9's quicksort did.

// --- 2. A heap's height is not a hope, it is a guarantee -----------
// A binary heap is a COMPLETE tree stored in an array: node i has
// children 2i+1 and 2i+2. Completeness is not maintained by clever
// rebalancing -- it is maintained by the array layout, and cannot
// fail. Hence height = floor(lg n), exactly, always.

// --- 3. A BST's height depends entirely on the insertion order ----
// Same keys, same code, two orders. This is Chapter 2's lesson --
// the nature of the input, not its size -- applied to a structure
// rather than to a sort.

// --- 4. The constant, averaged -------------------------------------
// Quoting only the leading constants (2.988 and 1.386) would make the
// measurement look wrong. Both have lower-order corrections that are
// still large at these sizes:
//   height     ~ 4.311 ln n - 1.953 ln ln n + O(1)   (Reed, 2003)
//   mean depth ~ 2 ln n - 3 + O(1)
// So the honest comparison is against the corrected forms, with the
// bare asymptote shown beside them to make the gap visible.

// --- 5. And what the height costs, in seconds ----------------------
// Searching a degenerate BST is a linear scan with pointer chasing,
// which is worse than a linear scan of an array: same number of
// steps, none of them predictable to the prefetcher.
\end{lstlisting}

\textbf{What to notice.} Part 4 prints three columns for each quantity:
measured, corrected prediction, and bare asymptote. That is deliberate. A
table with only the first and last would have read as a 20\% failure of a
well-established theorem.

\smallskip
\textbf{Then try to break it.} Insert the keys in \emph{bit-reversal} order ---
$0, 8, 4, 12, 2, 10, \ldots$ for $n = 16$ --- and measure the height. Predict
first. This order is neither sorted nor random, and it produces a
\emph{perfectly balanced} tree, height exactly $\lg n$: better than random. The
lesson is that ``not sorted'' and ``random'' are not the same assumption, and
only one of them is what the theorem requires.
\end{lab}

\begin{checkpoint}
\begin{tightnum}
  \item Explain why a heap cannot become unbalanced, in terms of its
        representation rather than its operations.
  \item The naive bound for \algname{Build-Heap} is $\bigO{n\lg n}$ and the
        true answer $\bigTh{n}$. Identify precisely which step of the naive
        argument is loose, and evaluate $\sum_{j\ge0} j/2^{j}$.
  \item A colleague's BST benchmark shows excellent performance. What single
        question would you ask about the test data \emph{(CLO-2)}?
\end{tightnum}
\end{checkpoint}

\begin{chaptersummary}
\begin{tight}
  \item A binary heap's height is exactly $\lfloor \lg n \rfloor$, guaranteed
        by the array layout rather than maintained by an algorithm. Measured
        exact at every size up to $10^{8}$.
  \item \algname{Build-Heap} is $\bigTh{n}$, not $\bigTh{n\lg n}$: most nodes
        are near the leaves, and $\sum_{j\ge0} j/2^{j} = 2$ converges. The
        naive bound is correct but loose --- the pattern \chref{amortised}
        generalises.
  \item A BST has no such guarantee. Its height is decided by the insertion
        order: measured $2.3\lg n$ for shuffled keys and exactly $n$ for sorted
        ones.
  \item Measured: 100{,}000 searches took 11.7 ms in a random tree and 7.6
        seconds in a degenerate one at $n = 32{,}000$ --- a factor of 650,
        of which only about 20 is step count and the rest unpredictable memory
        access.
  \item Measured height$/\lg n$ was 2.46 against an asymptote of 2.99; the
        $-1.953\ln\ln n$ correction accounts for the gap. Asymptotic constants
        are limits, and the discarded terms can dominate at every realistic
        size.
  \item A heap is a priority queue, not a search structure: $\bigTh{1}$ to the
        extremum and $\bigTh{n}$ to anything else.
\end{tight}
\end{chaptersummary}

\begin{reviewq}
  \item \textbf{(MCQ)} \algname{Build-Heap} on $n$ elements runs in:
        (a)~$\bigTh{\lg n}$; (b)~$\bigTh{n}$; (c)~$\bigTh{n\lg n}$;
        (d)~$\bigTh{n^{2}}$.
  \item \textbf{(MCQ)} Inserting $1,2,3,\ldots,n$ into an unbalanced BST gives
        height: (a)~$\lfloor\lg n\rfloor$; (b)~$\bigTh{\lg n}$;
        (c)~$\bigTh{\sqrt n}$; (d)~$n$.
  \item \textbf{(Short)} Explain why a heap needs no rebalancing while a BST
        does, in terms of what enforces the shape.
  \item \textbf{(Short)} Why is searching a heap $\bigTh{n}$ despite its height
        being $\lg n$?
  \item \textbf{(Short)} State what ``randomly built BST'' means, and give a key
        set that is random but produces a degenerate tree.
  \item \textbf{(Applied)} Prove that a heap on $n$ elements has at most
        $\lceil n/2^{j+1}\rceil$ nodes of height $j$, and use it to complete the
        $\bigTh{n}$ bound for \algname{Build-Heap}.
  \item \textbf{(Applied)} You must support insert, delete-min and lookup-by-key
        on up to $10^{7}$ records arriving in timestamp order. Choose a
        structure, justify it against the table in Section~12.4, and say which
        measurement from this chapter most influenced you.
\end{reviewq}
